% !TEX root = main.tex
We restate the ring-switching construction of Diamond and
Posen~\cite{eurocrypt/DiamondP26} in the vector and matrix notation
of this paper. This formulation supports the weighted tensor
queries arising from the reductions of Section~\ref{sec:blocks}.

Let \(N=m\kappa\) and \([\ExtensionField:\Field]=\kappa\).
We commit to the fixed-basis packing \(\bm A\ext{\bm\beta}\)
of a matrix \(\bm A\in\Field^{m\times\kappa}\), and reduce
each query to one MLE opening of this commitment.
The verifier evaluates the induced public query by a recursion
on small base-field matrices.
We then instantiate the underlying PCS from a base-field linear
code using Blaze's code-switching construction.

\subsubsection{Given a PCS over the Extension Field}

\medskip\noindent{\bf Commitment and Opening Relation}

Let \(m,\kappa\) be powers of two, with \(\kappa\geq2\), and write
\(\ell=\log_2m\).
Fix an \(\Field\)-basis
\(\ext{\bm\beta}=(\ext\beta_1,\ldots,\ext\beta_\kappa)^\top\)
of \(\ExtensionField\), with \(\ext\beta_1=1\), and let \(\pc\)
be a multilinear PCS for vectors in \(\ExtensionField^m\).
Extension-field elements are stored in this basis, so packing
and coordinate extraction only reorganize their base-field coordinates.
The schoolbook cost bounds below use a polynomial basis.

\paragraph{Setup and commitment.}
Run \(\pc.\Setup(1^\lambda,m)\) to obtain
\(\mathsf{pp}_{\pc}\), and include the fields, basis, and dimensions
\(N,m,\kappa\) in the public parameters.
Pack each row of the witness matrix into one extension-field element:
\[
    \bm A=\concat{\bm a^{(1)},\ldots,\bm a^{(\kappa)}}
        \in\Field^{m\times\kappa},
    \qquad
    \ext{\bm a}=\bm A\ext{\bm\beta}\in\ExtensionField^m.
\]
The commitment is
\(\mathsf{cm}=\pc.\Com(\mathsf{pp}_{\pc},\ext{\bm a};
\rho_{\pc})\), where \(\rho_{\pc}\) is the underlying commitment
randomness. The prover retains \(\bm A\) and \(\rho_{\pc}\).
Since the basis is fixed, binding to \(\ext{\bm a}\) binds the
prover to a unique matrix \(\bm A\in\Field^{m\times\kappa}\).

\paragraph{Opening relation.}
A query specifies column weights
\(\ext{\bm\rho}\in\ExtensionField^\kappa\), a tensor vector
\(\ext{\bm\chi}\in\ExtensionField^m\), and a claimed value
\(\ext u\in\ExtensionField\). The opening relation is:
\begin{equation}
    \ext u
    =\langle\bm A\ext{\bm\rho},\ext{\bm\chi}\rangle
    =\ext{\bm\rho}^{\top}\bm A^\top\ext{\bm\chi}.
\label{eq:pcs-opening-relation}
\end{equation}
The query is supplied by its two-entry tensor factors:
\[
    \ext{\bm\chi}
    =\bigotimes_{j=1}^{\ell}\ext{\bm\chi}^{(j)},
    \qquad
    \ext{\bm\chi}^{(j)}
    =
    \begin{pmatrix}
        \ext\chi^{(j)}_0\\
        \ext\chi^{(j)}_1
    \end{pmatrix}
    \in\ExtensionField^2.
\]
Ordinary MLE queries use the tensor spans of the corresponding
parts of the evaluation point for both \(\ext{\bm\rho}\) and
\(\ext{\bm\chi}\). Allowing arbitrary two-entry factors also
supports the coordinatewise products of tensor queries arising
in the weighted Hadamard reduction.

\medskip\noindent{\bf Packed Opening Protocol}

The commitment uses the fixed weights \(\ext{\bm\beta}\), whereas
the query uses \(\ext{\bm\rho}\). We first claim the column
evaluations, then use coordinate decomposition and a random
projection to reduce their correctness to an inner product with
the committed vector. The verifier evaluates the resulting public
query locally from its tensor factors.

\paragraph{Reduce the column claims.}
Write the public query and the claimed column values in base-field
coordinates:
\[
    \ext{\bm\chi}=\bm\Xi\ext{\bm\beta},
    \qquad
    \ext{\bm v}=\bm V\ext{\bm\beta},
    \qquad
    \bm\Xi\in\Field^{m\times\kappa},
    \quad
    \bm V\in\Field^{\kappa\times\kappa}.
\]
The honest prover computes \(\bm V=\bm A^\top\bm\Xi\) and sends
\(\ext{\bm v}=\bm V\ext{\bm\beta}\).
The verifier extracts \(\bm V\) from this message and checks
\(\ext u=\ext{\bm\rho}^{\top}\ext{\bm v}\).
The matrix \(\bm\Xi\) is defined implicitly by the public query.

Since \(\bm A\) has base-field entries, the remaining claim is
equivalent to a base-field matrix identity:
\[
    \ext{\bm v}=\bm A^\top\ext{\bm\chi}
    \quad\Longleftrightarrow\quad
    \bm V=\bm A^\top\bm\Xi.
\]
Packing its columns with the fixed basis introduces the committed
vector \(\ext{\bm a}\):
\begin{equation}
    \bm V=\bm A^\top\bm\Xi
    \quad\Longleftrightarrow\quad
    \ext{\bm\beta}^{\top}\bm V
    =\ext{\bm a}^{\top}\bm\Xi.
\label{eq:pcs-encoded-matrix}
\end{equation}
After receiving \(\ext{\bm v}\), the verifier samples
\(\ext{\bm r}\sample\ExtensionField^\kappa\), independently of
the original query weights, and projects
Eq.~\eqref{eq:pcs-encoded-matrix} to one inner-product claim:
\begin{equation}
    \ext c
    =\ext{\bm\beta}^{\top}\bm V\ext{\bm r}
    =\langle\ext{\bm a},\bm\Xi\ext{\bm r}\rangle.
\label{eq:pcs-random-opening}
\end{equation}
The verifier computes \(\ext c\) by interpreting the columns of
\(\bm V\) as extension-field elements and taking their inner
product with \(\ext{\bm r}\).

Apply Fact~\ref{fact:ip} over \(\ExtensionField\) to
Eq.~\eqref{eq:pcs-random-opening}.
Let \(\ext{\bm y}\in\ExtensionField^\ell\) be its terminal point
and set \(\ext{\bm\xi}=\tensorspan{\ext{\bm y}}\).
The prover's terminal values
\(\ext w_{\mathsf{com}},\ext w\in\ExtensionField\) give two claims:
\begin{equation}
    \ext w_{\mathsf{com}}
    =\langle\ext{\bm a},\ext{\bm\xi}\rangle,
    \qquad
    \ext w=\ext{\bm\xi}^{\top}\bm\Xi\ext{\bm r}.
\label{eq:pcs-terminal-claims}
\end{equation}
The first is an MLE opening of the committed vector.
The verifier checks the second by the following matrix recursion.

\paragraph{Evaluate the public query.}
Write \(\ext{\bm\xi}=\bm U\ext{\bm\beta}\), where
\(\bm U\in\Field^{m\times\kappa}\).
It suffices to compute the small matrix:
\begin{equation}
    \bm W=\bm U^\top\bm\Xi\in\Field^{\kappa\times\kappa},
    \qquad
    \ext w=\ext{\bm\beta}^{\top}\bm W\ext{\bm r}.
\label{eq:pcs-query-matrix}
\end{equation}
Both \(\bm U\) and \(\bm\Xi\) are determined by public tensor
factors. Their product can be evaluated without forming either
length-\(m\) coordinate table.

For \(\ext z\in\ExtensionField\), let
\(\bm M_{\ext z}\in\Field^{\kappa\times\kappa}\) be the coordinate
matrix of multiplication by \(\ext z\), characterized by:
\begin{equation}
    \ext{\bm\beta}^{\top}\bm M_{\ext z}
    =\ext z\,\ext{\bm\beta}^{\top}.
\label{eq:pcs-multiplication-matrix}
\end{equation}
Thus left multiplication by \(\bm M_{\ext z}\) scales each
coordinate column by \(\ext z\) in \(\ExtensionField\), and right
multiplication by \(\bm M_{\ext z}^\top\) scales each coordinate row.
In particular, \(\bm M_{1-\ext z}=\bm I-\bm M_{\ext z}\).

To expose the recursion, let \(\bm U_j,\bm\Xi_j\) denote the
coordinate matrices for the suffix factors \(j,\ldots,\ell\):
\[
    \bm U_j\ext{\bm\beta}
    =\bigotimes_{i=j}^{\ell}
        \begin{pmatrix}1-\ext y_i\\ \ext y_i\end{pmatrix},
    \qquad
    \bm\Xi_j\ext{\bm\beta}
    =\bigotimes_{i=j}^{\ell}\ext{\bm\chi}^{(i)}.
\]
Here \(\bm U_1=\bm U\), \(\bm\Xi_1=\bm\Xi\), and
\(\bm U_{\ell+1}=\bm\Xi_{\ell+1}=\bm e_1^\top\), where
\(\bm e_1=(1,0,\ldots,0)^\top\in\Field^\kappa\).
The two possible values of the first suffix bit give the block
decompositions:
\begin{equation}
\begin{aligned}
    \bm U_j
    &=
    \begin{bmatrix}
        \bm U_{j+1}\bm M_{1-\ext y_j}^\top\\
        \bm U_{j+1}\bm M_{\ext y_j}^\top
    \end{bmatrix},\\
    \bm\Xi_j
    &=
    \begin{bmatrix}
        \bm\Xi_{j+1}\bm M_{\ext\chi^{(j)}_0}^\top\\
        \bm\Xi_{j+1}\bm M_{\ext\chi^{(j)}_1}^\top
    \end{bmatrix}.
\end{aligned}
\label{eq:pcs-query-blocks}
\end{equation}
Consequently, the verifier computes
\(\bm W_j=\bm U_j^\top\bm\Xi_j\) by initializing
\(\bm W_{\ell+1}=\bm e_1\bm e_1^\top\) and processing
\(j=\ell,\ldots,1\):
\begin{equation}
\begin{aligned}
    \bm W_j
    &=
    \bm M_{1-\ext y_j}\bm W_{j+1}
        \bm M_{\ext\chi^{(j)}_0}^\top
    +\bm M_{\ext y_j}\bm W_{j+1}
        \bm M_{\ext\chi^{(j)}_1}^\top\\
    &=
    \bm W_{j+1}\bm M_{\ext\chi^{(j)}_0}^\top
    +\bm M_{\ext y_j}\bm W_{j+1}
        \bm M_{\ext\chi^{(j)}_1-\ext\chi^{(j)}_0}^\top.
\end{aligned}
\label{eq:pcs-query-recurrence}
\end{equation}
The output is \(\bm W=\bm W_1\).
The verifier never forms the suffix coordinate matrices.
Each step of the second line uses at most \(3\kappa\)
extension-field multiplications, implemented by row and column
scaling without materializing \(\bm M_{\ext z}\).
The recursion takes \(O(\kappa\ell)\) extension-field
multiplications and \(O(\kappa^2\ell)\) base-field additions and
coordinate moves, using \(O(\kappa^2)\) base-field coordinates
of working memory. It applies in every characteristic.

The verifier packs the columns of \(\bm W\), computes
\(\ext{\bm\beta}^{\top}\bm W\ext{\bm r}\), and checks the second
identity in Eq.~\eqref{eq:pcs-query-matrix}.
This computation is exact and requires no additional prover message.

\paragraph{Terminal opening.}
Finally, the parties authenticate \(\ext w_{\mathsf{com}}\)
at \(\ext{\bm y}\) through the underlying PCS:
\[
    \pc.\Open\!\left(
        \mathsf{pp}_{\pc},\mathsf{cm},
        \ext w_{\mathsf{com}},\ext{\bm y}
        :\ext{\bm a},\rho_{\pc}
    \right).
\]
The verifier accepts if the scalar checks, the inner-product
reduction, and the underlying opening accept.
Algorithm~\ref{alg:packed-pcs} in Appendix~\ref{appx:pseudocode}
summarizes the protocol.

\medskip\noindent{\bf Security and Complexity}

\begin{proposition}[Packed polynomial commitment scheme]
\label{proposition:packed-pcs}
Let \(\pc\) bind the prover to a vector in \(\ExtensionField^m\),
with evaluation-soundness error \(\epsilon_{\mathsf{PCS}}\) for
adaptively chosen MLE claims about that vector.
The packed commitment is binding and the opening protocol is
complete, with evaluation-soundness error at most:
\begin{equation}
    \epsilon_{\mathsf{PCS}}
    +\frac{1+2\ell}{|\ExtensionField|}.
\label{eq:pcs-soundness}
\end{equation}
The opening reduction requires no additional lower bound on
\(|\Field|\).

Excluding the underlying PCS opening, the prover's arithmetic cost is:
\begin{equation}
    O(\kappa N)\,\bb
    +O(N)\,\be
    +O(m+\kappa)\,\ee.
\label{eq:pcs-prover-cost}
\end{equation}
The verifier's arithmetic cost is:
\begin{equation}
    O\bigl(\kappa^2(1+\ell)\bigr)\,\bb
    +O\bigl(\kappa(1+\ell)\bigr)\,\ee.
\label{eq:pcs-verifier-cost}
\end{equation}
The verifier bound includes base-field additions and coordinate
reorganization in the matrix recursion.
When \(m\geq\kappa\), schoolbook arithmetic gives
\(O(\kappa N)\) base-field operations for the prover.
Peak prover memory is \(O(N+\kappa^2)\) base-field coordinates,
or \(O(N)\) when \(m\geq\kappa\), excluding the workspace of the
underlying PCS.

With \(\ext u\) already part of the public statement,
prover-to-verifier communication, excluding the underlying
opening, consists of \(\kappa+2\ell+2\) extension-field elements.
The packed commitment inherits hiding from the underlying
commitment. The opening protocol reveals the column evaluations
\(\ext{\bm v}\), so this inheritance does not establish
zero knowledge of the opening transcript.
\end{proposition}

The proof is deferred to Appendix~\ref{appx:proof-pcs}.

\subsubsection{Given a Linear Code over the Base Field}

A linear code over \(\Field\) also encodes packed messages over
\(\ExtensionField\): applying its generator matrix commutes with
basis packing. We use this observation to instantiate the preceding
construction via Blaze's code-switching
protocol~\cite{eurocrypt/BrehmCFRRZ25}, with a BaseFold
code~\cite{crypto/ZeilbergerCF24} as the target. We use Blaze under
its source-encoding and proximity assumptions.

Let \(\Code:\Field^m\to\Field^n\) be a linear code with
generator matrix \(\bm G\in\Field^{n\times m}\).
Its \(\kappa\)-way interleaving encodes the columns of
\(\bm A\in\Field^{m\times\kappa}\) independently.
Using the same fixed basis,
the resulting codeword has the equivalent packed representation:
\begin{equation}
\begin{aligned}
    \bm C&=\bm G\bm A\in\Field^{n\times\kappa},\\
    \ext{\bm c}&=\bm C\ext{\bm\beta}
       =\bm G(\bm A\ext{\bm\beta})
       =\Code_{\ExtensionField}(\ext{\bm a}).
\end{aligned}
\label{eq:pcs-code-packing}
\end{equation}
Here \(\Code_{\ExtensionField}\) is the scalar extension defined
by the same generator matrix; it has the same rate and relative
distance as \(\Code\).
Thus a row of \(\bm C\) represents one extension-field codeword
symbol. Committing to these rows supplies Blaze's source codeword
oracle for \(\ext{\bm a}=\bm A\ext{\bm\beta}\).

Apply Blaze's code-switching construction over \(\ExtensionField\)
to this oracle, switching to a BaseFold encoding of the same
length-\(m\) message \(\ext{\bm a}\).
The resulting MLE openings instantiate the underlying PCS \(\pc\)
in the preceding construction.
The packed opening reduction handles the challenge-dependent weights
\(\ext{\bm\rho}\), and code switching authenticates its terminal
fixed-basis opening. We include the costs and soundness error of
the chosen Blaze instantiation, including source-encoding
consistency, in the underlying PCS costs and
\(\epsilon_{\mathsf{PCS}}\), respectively.

If \(\Code\) is constant-rate and linear-time encodable, forming
\(\bm C\) costs \(O(N)\) base-field operations and its payload
occupies \(O(N\log_2|\Field|)\) bits, excluding Merkle-tree hashes.
The BaseFold stage works on \(m=N/\kappa\) extension-field message
symbols. Thus the source encoding retains the witness's small-field
representation, and the extension-field opening operates on the
shorter packed message.
